% Document class and font size
\documentclass[a4paper,9pt]{extarticle}

% Packages
\usepackage[utf8]{inputenc} % For input encoding
\usepackage{geometry} % For page margins
\geometry{letterpaper, margin=0.75in} % Set paper size and margins
\usepackage{titlesec} % For section title formatting
\usepackage{enumitem} % For itemized list formatting
\usepackage{hyperref} % For hyperlinks
\usepackage{tabularx}
\usepackage{fancyhdr}

% Formatting
\setlist{noitemsep} % Removes item separation
\titleformat{\section}{\large\bfseries}{\thesection}{1em}{}[\titlerule] % Section title format
\titlespacing*{\section}{0pt}{\baselineskip}{\baselineskip} % Section title spacing

% Begin document
\begin{document}

% Disable page numbers
\pagestyle{fancy}
\renewcommand{\headrulewidth}{0pt}
\fancyhead{}
\fancyhead[L]{\textit{Anrui Wang}}
\fancyhead[R]{\textit{\# Month \#Year}}
\thispagestyle{empty} % Remove header from the first page

% Header
\begin{flushleft}
\textbf{\LARGE Anrui Wang}\\[2pt] % Name
CS student
\end{flushleft}

% PLACEHOLDER-APPEND-MARKER

\section*{RESEARCH INTERESTS}
\noindent
AI for Science, Single-cell Representation Learning, Graph Representation Learning

\section*{RESEARCH EXPERIENCE}
\noindent
\textbf{STATE-Based Multi-Omics Modeling for Synthetic-Lethality Prediction} \\
\textit{Supervised by Prof. Jie Zheng} \hfill Nov. 2025 -- Present
\begin{itemize}
    \item Framed and independently led a staged research program toward synthetic-lethal (SL) target discovery, predicting gene-pair SL interactions for cancer-selective drug-target prioritization.
    \item Integrated multi-omics evidence into SL pair prediction: Perturb-seq perturbation transcriptomics, DepMap functional dependency (CRISPR GeneEffect), together with ESM2 protein-sequence embeddings.
    \item Modeled perturbation responses as cell-state representations predictive of cancer dependency, developing from variational single-cell representations (scVI) and multiple-instance pooling to a prototype-based distribution-regression model that recovers dependency signal missed by pseudobulk and attention baselines.
    \item Built cross-cell-line selectivity features for context-specific dependency, and evaluated SL prediction on a gene-pair benchmark with adversarial leakage.
\end{itemize}

\noindent
\textbf{Topology-Constrained Deep Learning for Inductive Interactome Reconstruction} \\
\textit{Supervised by Prof. Xuming He} \hfill Jan. 2026 -- Present
\begin{itemize}
    \item Reframed protein--protein interaction prediction from independent pairwise classification into inductive reconstruction of the interaction network.
    \item Designed a topology-constrained conditional generator modeling the full adjacency structure of an unseen protein set via a decomposed edge formulation (pairwise compatibility, hub propensity, community membership, set-level density).
    \item Instilled topological realism with a training-only topology teacher (masked-edge reconstruction distillation) and differentiable graph-level objectives, while keeping the deployed model strictly inductive on intrinsic protein features (ESM).
    \item Established credible baselines through systematic architecture ablations, and reproduction of SOTA models.
\end{itemize}

\section*{SKILLS}
\begin{itemize}
    \item \textbf{Programming:} Python, PyTorch
    \item \textbf{ML / Tooling:} HuggingFace Accelerate, single-cell foundation models, ESM series, graph neural networks
    \item \textbf{Coding Agents:} highly skilled with agentic coding tools, including Claude Code and Codex
    \item \textbf{Domains:} single-cell genomics, cancer dependency (DepMap), synthetic lethality, protein--protein interaction networks
\end{itemize}

\section*{SELECTED COURSES}
\noindent
Machine Learning, Deep Learning, AI for Science and Engineering, Bioinformatics

\end{document}
